% fig1_overview_v2 — 2-row 7-panel narrative overview
% Canvas: 14cm × 9cm  →  \resizebox{178mm}{!}  (target NatComm 178mm × ~120mm)
% Row 1 (y=5.00..9.00): A (x=0..3.50) | B (x=3.55..6.85) | C HERO (x=6.90..13.95, ~1.5×)
% Inter-row gap (y=4.50..5.00) carries the C → Row 2 "convergent evidence" arrow.
% Row 2 (y=0.10..4.45): D (x=0..4.45) | E (x=4.50..6.95) | F (x=7.00..10.45) | G (x=10.50..13.95)
%
% Implementation status (design.md v2.1 §10):
%   Panel C   — IMPLEMENTED (HERO #1)
%   Panels A, B, D, E, F, G — placeholder regions (filled in subsequent rounds)
%
% Spring-layout / chemfig packages intentionally NOT imported here:
% Panel A draft used graphdrawing/force which design.md §7 forbids; chemfig
% will be re-introduced when Panel A's S₈ corner inset is authored.

\documentclass[border=4pt]{standalone}
\usepackage{polymer-paper-preamble}

\begin{document}
\resizebox{178mm}{!}{%
\begin{tikzpicture}[
  every node/.style={font=\sffamily\fontsize{8}{10}\selectfont},
  interArrow/.style={-{Stealth[length=7pt,width=5pt]}, cGray!55!black,
                     line width=1.4pt},
  % --- 3-tier inline label hierarchy (design.md §6, fig2 calibration) ---
  labelStrong/.style={font=\sffamily\bfseries\fontsize{8.5}{10.2}\selectfont},
  labelStd/.style={font=\sffamily\fontsize{7.5}{9}\selectfont},
  labelMute/.style={font=\sffamily\itshape\fontsize{7}{8.4}\selectfont,
                    text=cGray!80!black},
]

% ============================================================
% Figure-wide background wash (single canvas tone — §2.2 compliant)
% ============================================================
\fill[cLGray!8, rounded corners=3mm] (-0.10, -0.10) rectangle (14.10, 9.10);

% ============================================================
% Panel scaffold reminders — text-only at low opacity (§2.2 compliant)
% Bbox refs: A x=0..3.50 y=5..9 | B x=3.55..6.85 y=5..9
%            D x=0..4.45 y=0.1..4.45 | E x=4.50..6.95 | F x=7..10.45 | G x=10.5..13.95
% ============================================================

% ============================================================
% Panel A — Sulfur-rich polymer (chemistry-accurate identifier) — design.md §4 v2.2
% Spec: 2 DIB benzene rings linked by polysulfide bridge chain + dangling chains
% at 1,3-meta vertices toward panel edges (suggesting larger network).
% + S₈ ring inset upper-right + transformation arrow.
% Bezier chains (sine forbidden). Manual hexagonal aromatic notation.
% Forbidden: 3+ DIB rings, S₈ in center, generic atom-bond mesh.
% Connection out: 우측 화살표 clean → Panel B.
% ============================================================

% --- Subtle amber wash (figure-wide consistency, centered on bridge midpoint) ---
\fill[cAmber!08] (1.75, 6.95) circle (1.55cm);

% --- DIB benzene ring helper (r=0.42, hexagonal aromatic notation) ---
\newcommand{\dibRing}{
  \pgfmathsetmacro{\dibR}{0.42}
  \draw[cGray!85!black, line width=0.70pt]
    (0:\dibR) -- (60:\dibR) -- (120:\dibR) -- (180:\dibR) --
    (240:\dibR) -- (300:\dibR) -- cycle;
  \draw[cGray!85!black, line width=0.55pt] (66:\dibR*0.78) -- (114:\dibR*0.78);
  \draw[cGray!85!black, line width=0.55pt] (186:\dibR*0.78) -- (234:\dibR*0.78);
  \draw[cGray!85!black, line width=0.55pt] (306:\dibR*0.78) -- (354:\dibR*0.78);
}

% Two-DIB centers (meta-substituted, bridged):
% DIB1 (upper-left)  center = (1.10, 7.60), r = 0.42
%   vertex @180° (toward panel-left dangling chain): (0.680, 7.600)
%   vertex @300° (toward bridge / DIB2):             (1.310, 7.236)
% DIB2 (lower-right) center = (2.40, 6.30), r = 0.42
%   vertex @120° (toward bridge / DIB1):             (2.190, 6.664)
%   vertex   @0° (toward panel-right dangling chain):(2.820, 6.300)
%
% Isopropenyl linker length = 0.13cm (short straight gray, ring vertex → chain anchor):
%   DIB1 @180° → (0.550, 7.600)
%   DIB1 @300° → (1.375, 7.123)
%   DIB2 @120° → (2.125, 6.777)
%   DIB2 @0°   → (2.950, 6.300)

% --- Isopropenyl linker bonds (4 short gray straight) ---
\draw[cGray!75!black, line width=0.80pt] (0.680, 7.600) -- (0.550, 7.600);
\draw[cGray!75!black, line width=0.80pt] (1.310, 7.236) -- (1.375, 7.123);
\draw[cGray!75!black, line width=0.80pt] (2.190, 6.664) -- (2.125, 6.777);
\draw[cGray!75!black, line width=0.80pt] (2.820, 6.300) -- (2.950, 6.300);

% --- Bridge chain (DIB1 ↔ DIB2): polysulfide -S-S- linkage carrying network identity ---
\draw[cAmber!85!black, line width=0.95pt]
  (1.375, 7.123)
  .. controls (1.55, 7.10) and (1.65, 6.78) .. (1.80, 6.85)
  .. controls (1.95, 6.92) and (2.00, 6.65) .. (2.125, 6.777);
\node[font=\sffamily\fontsize{6.5}{7.8}\selectfont, text=cAmber!85!black,
      anchor=west, fill=white, inner sep=0.6pt]
  at (1.82, 7.05) {S\,--\,S};

% --- Dangling chain 1 (DIB1 @180° → panel-left edge, fade-out) ---
\draw[cAmber!85!black, line width=0.95pt]
  (0.550, 7.600)
  .. controls (0.40, 7.65) and (0.30, 7.85) .. (0.18, 7.92);
\draw[cAmber!85!black, line width=0.95pt, opacity=0.32]
  (0.18, 7.92)
  .. controls (0.05, 7.99) and (-0.05, 8.05) .. (-0.10, 8.10);
\node[font=\sffamily\fontsize{6.5}{7.8}\selectfont, text=cAmber!85!black]
  at (0.30, 8.10) {S\,--\,S\,--\,S};

% --- Dangling chain 2 (DIB2 @0° → panel-right edge, fade-out) ---
\draw[cAmber!85!black, line width=0.95pt]
  (2.950, 6.300)
  .. controls (3.10, 6.20) and (3.22, 6.05) .. (3.32, 5.92);
\draw[cAmber!85!black, line width=0.95pt, opacity=0.32]
  (3.32, 5.92)
  .. controls (3.42, 5.83) and (3.50, 5.78) .. (3.55, 5.72);
\node[font=\sffamily\fontsize{6.5}{7.8}\selectfont, text=cAmber!85!black]
  at (3.18, 6.20) {S\,--\,S};

% --- DIB rings drawn LAST so they cover linker endpoints cleanly ---
\begin{scope}[shift={(1.10, 7.60)}] \dibRing \end{scope}
\begin{scope}[shift={(2.40, 6.30)}] \dibRing \end{scope}

% --- S₈ ring inset (upper-right corner, manual octagon, ~15-20% panel area) ---
\begin{scope}[shift={(2.95, 8.50)}]
  \pgfmathsetmacro{\sR}{0.32}
  \foreach \i in {0,...,7} {
    \pgfmathsetmacro{\sang}{90 - \i*45}
    \pgfmathsetmacro{\snext}{90 - (\i+1)*45}
    \draw[cAmber!75!black, line width=0.55pt]
      ({cos(\sang)*\sR}, {sin(\sang)*\sR}) --
      ({cos(\snext)*\sR}, {sin(\snext)*\sR});
  }
  \foreach \i in {0,...,7} {
    \pgfmathsetmacro{\sang}{90 - \i*45}
    \node[font=\sffamily\bfseries\fontsize{5}{6}\selectfont,
          text=cAmber!85!black, fill=white, inner sep=0.3pt]
      at ({cos(\sang)*\sR}, {sin(\sang)*\sR}) {S};
  }
\end{scope}

% --- Transformation arrow: S₈ → DIB (inverse vulcanization 서사) ---
\draw[cAmber!85!black, dashed, line width=0.85pt,
      -{Stealth[length=4.5pt,width=3.5pt]}]
  (2.78, 8.20) to[out=-115, in=40] (2.10, 7.30);

% --- "S₈" label below ring ---
\node[labelStd, text=cAmber!75!black, anchor=north]
  at (2.95, 8.10) {S$_8$};

% --- Caption ---
\node[font=\sffamily\bfseries\fontsize{7.5}{9}\selectfont,
      text=cAmber!75!black] at (1.55, 5.42) {Sulfur-rich polymer};
\node[labelMute]            at (1.55, 5.17) {DIB-linked polysulfide network};

% --- A → B connector (clean) ---
\draw[cGray!55!black, line width=0.55pt,
      -{Stealth[length=4pt,width=2.5pt]}]
  (3.45, 6.80) -- (3.62, 6.80);

% ============================================================
% Panel B — S₆₀..S₈₅ chain length (soft tunability)  x: 3.55..6.85, y: 5..9
% 4 horizontal atom-explicit chains (subitizing limit), top→bottom short→long.
% Each chain: gold S atoms + gray cross-link C atoms alternating along zigzag
% backbone (design.md §4: "gold S 원자 + gray cross-link 원자 + bond").
% Labels (S₆₀, S₈₅) anchored EAST of chain start to avoid overlap with terminal atom.
% Atom counts: S60=4S+3C(7), mid=5S+4C(9)/6S+5C(11), S85=8S+7C(15) → 2.1× length ratio.
% Bond pitch=0.20cm, zigzag amplitude=0.10cm, chain start x=3.95 (uniform).
% ============================================================

% Helper: one chain row — backbone polyline + S atoms (even idx) + C atoms (odd idx)
% Args: #1=y_center, #2=n_atoms (must be odd: 7, 9, 11, 15)

% --- Chain 1 (top, S₆₀, n=7: 4 S + 3 C) ---
\node[labelStrong, text=cAmber!85!black, anchor=east] at (3.85, 8.40) {S$_{60}$};
\draw[cAmber!75!black, line width=0.70pt]
  (3.95, 8.50) -- (4.15, 8.30) -- (4.35, 8.50) -- (4.55, 8.30)
  -- (4.75, 8.50) -- (4.95, 8.30) -- (5.15, 8.50);
\foreach \sx/\sy in {3.95/8.50, 4.35/8.50, 4.75/8.50, 5.15/8.50} {
  \node[circle, fill=cAmber!85!black, draw=cAmber!95!black,
        inner sep=0pt, minimum size=1.3mm, line width=0.25pt] at (\sx, \sy) {};
}
\foreach \cx/\cy in {4.15/8.30, 4.55/8.30, 4.95/8.30} {
  \node[circle, fill=cGray!65, draw=cGray!85!black,
        inner sep=0pt, minimum size=0.85mm, line width=0.20pt] at (\cx, \cy) {};
}

% --- Chain 2 (n=9: 5 S + 4 C) ---
\draw[cAmber!75!black, line width=0.70pt]
  (3.95, 7.75) -- (4.15, 7.55) -- (4.35, 7.75) -- (4.55, 7.55)
  -- (4.75, 7.75) -- (4.95, 7.55) -- (5.15, 7.75) -- (5.35, 7.55) -- (5.55, 7.75);
\foreach \sx/\sy in {3.95/7.75, 4.35/7.75, 4.75/7.75, 5.15/7.75, 5.55/7.75} {
  \node[circle, fill=cAmber!85!black, draw=cAmber!95!black,
        inner sep=0pt, minimum size=1.3mm, line width=0.25pt] at (\sx, \sy) {};
}
\foreach \cx/\cy in {4.15/7.55, 4.55/7.55, 4.95/7.55, 5.35/7.55} {
  \node[circle, fill=cGray!65, draw=cGray!85!black,
        inner sep=0pt, minimum size=0.85mm, line width=0.20pt] at (\cx, \cy) {};
}

% --- Chain 3 (n=11: 6 S + 5 C) ---
\draw[cAmber!75!black, line width=0.70pt]
  (3.95, 7.00) -- (4.15, 6.80) -- (4.35, 7.00) -- (4.55, 6.80)
  -- (4.75, 7.00) -- (4.95, 6.80) -- (5.15, 7.00) -- (5.35, 6.80)
  -- (5.55, 7.00) -- (5.75, 6.80) -- (5.95, 7.00);
\foreach \sx/\sy in {3.95/7.00, 4.35/7.00, 4.75/7.00, 5.15/7.00, 5.55/7.00, 5.95/7.00} {
  \node[circle, fill=cAmber!85!black, draw=cAmber!95!black,
        inner sep=0pt, minimum size=1.3mm, line width=0.25pt] at (\sx, \sy) {};
}
\foreach \cx/\cy in {4.15/6.80, 4.55/6.80, 4.95/6.80, 5.35/6.80, 5.75/6.80} {
  \node[circle, fill=cGray!65, draw=cGray!85!black,
        inner sep=0pt, minimum size=0.85mm, line width=0.20pt] at (\cx, \cy) {};
}

% --- Chain 4 (bottom, S₈₅, n=15: 8 S + 7 C) ---
\node[labelStrong, text=cAmber!85!black, anchor=east] at (3.85, 6.25) {S$_{85}$};
\draw[cAmber!75!black, line width=0.70pt]
  (3.95, 6.25) -- (4.15, 6.05) -- (4.35, 6.25) -- (4.55, 6.05)
  -- (4.75, 6.25) -- (4.95, 6.05) -- (5.15, 6.25) -- (5.35, 6.05)
  -- (5.55, 6.25) -- (5.75, 6.05) -- (5.95, 6.25) -- (6.15, 6.05)
  -- (6.35, 6.25) -- (6.55, 6.05) -- (6.75, 6.25);
\foreach \sx/\sy in {3.95/6.25, 4.35/6.25, 4.75/6.25, 5.15/6.25, 5.55/6.25,
                     5.95/6.25, 6.35/6.25, 6.75/6.25} {
  \node[circle, fill=cAmber!85!black, draw=cAmber!95!black,
        inner sep=0pt, minimum size=1.3mm, line width=0.25pt] at (\sx, \sy) {};
}
\foreach \cx/\cy in {4.15/6.05, 4.55/6.05, 4.95/6.05, 5.35/6.05, 5.75/6.05,
                     6.15/6.05, 6.55/6.05} {
  \node[circle, fill=cGray!65, draw=cGray!85!black,
        inner sep=0pt, minimum size=0.85mm, line width=0.20pt] at (\cx, \cy) {};
}

% --- Bottom axis arrow + label ---
\draw[cGray!65!black, line width=0.45pt,
      -{Stealth[length=3.5pt,width=2.5pt]}]
  (3.95, 5.55) -- (6.75, 5.55);
\node[labelMute] at (5.35, 5.32) {S-chain length};

% --- B → C connector (clean) ---
\draw[cGray!55!black, line width=0.55pt,
      -{Stealth[length=4pt,width=2.5pt]}]
  (6.78, 6.85) -- (6.93, 6.85);

% ============================================================
% Panel C — Localized traps  (HERO #1)   x: 6.90..13.95, y: 5.00..9.00
% ============================================================

% --- Faded polymer network background (4 chains + cross-links, opacity 0.42) ---
% Recovers design.md §4 "rich amorphous network" intent — 4 chains stacked
% vertically with 4 cross-link bonds, ~20 S atoms total. Atom size bumped to
% 0.55mm so polymer reads through the fade. opacity raised 0.28→0.42 per
% critique C002 (backdrop too faint, region read as 4-panel split).
\begin{scope}[opacity=0.42]
  % Chain 1 — top-left (highest)
  \draw[cGray!75!black, line width=0.80pt, smooth, samples=80,
        domain=7.40:10.40]
        plot (\x, {8.20 + 0.17*sin( (\x-7.40)*600 ) });
  \foreach \sx in {7.55, 8.15, 8.75, 9.35, 9.95} {
    \fill[cAmber!80!black] (\sx, 8.37) circle (0.55mm);
  }

  % Chain 2 — top-right (highest)
  \draw[cGray!75!black, line width=0.80pt, smooth, samples=80,
        domain=10.80:13.70]
        plot (\x, {8.05 + 0.17*sin( (\x-10.80)*600 ) });
  \foreach \sx in {10.95, 11.55, 12.15, 12.75, 13.35} {
    \fill[cAmber!80!black] (\sx, 8.22) circle (0.55mm);
  }

  % Chain 3 — mid-left (lower row, slight x offset for amorphous feel)
  \draw[cGray!75!black, line width=0.80pt, smooth, samples=70,
        domain=7.50:10.30]
        plot (\x, {7.40 + 0.17*sin( (\x-7.50)*600 ) });
  \foreach \sx in {7.65, 8.25, 8.85, 9.45, 10.05} {
    \fill[cAmber!80!black] (\sx, 7.57) circle (0.55mm);
  }

  % Chain 4 — mid-right (lower row)
  \draw[cGray!75!black, line width=0.80pt, smooth, samples=70,
        domain=10.85:13.60]
        plot (\x, {7.30 + 0.17*sin( (\x-10.85)*600 ) });
  \foreach \sx in {11.00, 11.60, 12.20, 12.80, 13.40} {
    \fill[cAmber!80!black] (\sx, 7.47) circle (0.55mm);
  }

  % Cross-link bonds (vertical, top↔mid chains — sparse, ~1 per quadrant)
  \draw[cGray!75!black, line width=0.55pt] (8.15, 8.37) -- (8.25, 7.57);
  \draw[cGray!75!black, line width=0.55pt] (9.35, 8.37) -- (9.45, 7.57);
  \draw[cGray!75!black, line width=0.55pt] (11.55, 8.22) -- (11.60, 7.47);
  \draw[cGray!75!black, line width=0.55pt] (12.75, 8.22) -- (12.80, 7.47);
\end{scope}

% --- Energy baseline reference ---
\draw[cGray!45, dashed, line width=0.30pt] (7.20, 6.85) -- (13.70, 6.85);

% --- Deep energy landscape — 3 wells, blue, amplitude 0.60 (left, briefing §3 Panel C) ---
% C001 BLOCKER fix: amplitudes swapped so left=deep (blue, briefing §6 deeper wells = blue).
\draw[cBlue!75!black, line width=1.30pt, smooth, samples=140,
      domain=7.30:9.80]
      plot (\x, {6.85 - 0.60*abs(sin( ((\x-7.30)/0.8333)*180 ))});

% --- Shallow energy landscape — 4 wells, red, amplitude 0.22 (right) ---
\draw[cRed!80!black, line width=1.00pt, smooth, samples=180,
      domain=10.05:13.65]
      plot (\x, {6.85 - 0.22*abs(sin( ((\x-10.05)/0.90)*180 ))});

% --- Deep electrons (3 × blue ⊖ at deep well bottoms, y=6.25) ---
\foreach \ex in {7.72, 8.55, 9.38} {
  \node[circle, draw=cBlue!85!black, fill=cBlue!55,
        inner sep=0pt, minimum size=1.7mm, line width=0.4pt,
        font=\sffamily\fontsize{5.5}{7}\selectfont, text=white]
    at (\ex, 6.25) {$-$};
}

% --- Shallow electrons (4 × red ⊖ at shallow well bottoms, y=6.63) ---
\foreach \ex in {10.50, 11.40, 12.30, 13.20} {
  \node[circle, draw=cRed!85!black, fill=cRed!55,
        inner sep=0pt, minimum size=1.7mm, line width=0.4pt,
        font=\sffamily\fontsize{5.5}{7}\selectfont, text=white]
    at (\ex, 6.63) {$-$};
}

% --- Region labels (inline only, labelStrong tier) ---
% C001 fix: left=deep (blue, briefing convention), right=shallow (red).
\node[labelStrong, text=cBlue!75!black] at (8.55, 5.85) {deep};
\node[labelStrong, text=cRed!80!black]  at (11.85, 5.85) {shallow};

% --- Subtitle (labelMute tier) ---
\node[labelMute] at (10.42, 5.32) {localized traps};

% ============================================================
% Panel D — Distributed release sequence + kinetic evidence   x: 0.00..4.45, y: 0.10..4.45
% Top: t1/t2 deep-well retention (blue) → t3/t4 shallow release (red).
% Bottom: multi-slope log-log I(t) plot with Debye reference.
% Concept labels only — composition-specific (S60/S85) forbidden by design.md.
% ============================================================

% --- Panel D title ---
\node[labelStrong, text=cGray!85!black] at (2.25, 4.28) {distributed release};

% --- t1..t4 de-trapping snapshots (briefing §3 Panel D) ---
\foreach \x/\lab/\col/\released in {
  0.75/t$_1$/cBlue/0,
  1.45/t$_2$/cBlue/0,
  2.65/t$_3$/cRed/1,
  3.35/t$_4$/cRed/1
} {
  \draw[cGray!70!black, line width=0.55pt, smooth]
    plot coordinates {(\x-0.33,3.42) (\x-0.16,3.24) (\x,3.16)
                      (\x+0.16,3.24) (\x+0.33,3.42)};
  \node[font=\sffamily\bfseries\fontsize{6.5}{7.8}\selectfont,
        text=\col!85!black] at (\x-0.28,3.84) {\lab};
  \node[circle, draw=\col!85!black, fill=\col!18,
        inner sep=0pt, minimum size=1.05mm, line width=0.28pt,
        font=\sffamily\fontsize{4}{4.8}\selectfont, text=\col!85!black]
        at (\x,3.22) {$-$};
  \ifnum\released=1
    \draw[\col!85!black, dashed, line width=0.45pt,
          -{Stealth[length=2.4pt,width=1.7pt]}]
      (\x+0.08,3.33) -- (\x+0.22,3.74);
  \fi
}
\node[labelMute] at (2.05, 3.34) {$\cdots$};
\draw[cGray!55!black, line width=0.45pt, -{Stealth[length=3pt,width=2pt]}]
  (0.55, 2.98) -- (3.68, 2.98);

% --- Plot axes ---
\draw[cGray, line width=0.40pt, -{Stealth[length=2.5pt,width=1.8pt]}]
  (0.40, 0.45) -- (4.30, 0.45);
\draw[cGray, line width=0.40pt, -{Stealth[length=2.5pt,width=1.8pt]}]
  (0.40, 0.45) -- (0.40, 2.35);

% --- non-Debye tail highlight (faint red wash, lower-right region) ---
% Wash alone carries the "deep-rich persists where Debye fails" message;
% explicit "non-Debye" text dropped to reduce label competition.
\fill[cRed!8] (3.00, 0.48) rectangle (4.28, 1.38);

% --- Debye reference (dashed gray, exponential → fast drop in log-log) ---
\draw[cGray!75!black, line width=0.65pt, dashed, smooth]
  plot coordinates {(0.55, 2.25) (1.50, 2.15) (2.20, 1.95)
                    (2.70, 1.55) (3.00, 0.92) (3.15, 0.52)};

% --- Shallow-rich I(t) — steep slope (high n), blue ---
\draw[cBlue!75!black, line width=0.95pt]
  (0.55, 2.55) -- (3.70, 0.65);

% --- Deep-rich I(t) — gentle slope (low n), red, persists in tail ---
% C007 fix: extend with shallower slope so non-Debye tail is unambiguous
% (deep-rich line clearly above Debye reference at long t).
\draw[cRed!80!black, line width=1.05pt]
  (0.55, 2.15) -- (4.25, 1.45);

% --- Curve callouts (concept labels — NO composition tags) ---
\node[labelStd, text=cBlue!75!black] at (1.75, 0.98) {shallow-rich};
\node[labelStd, text=cRed!80!black]  at (3.55, 1.60) {deep-rich};

% --- Main equation label (labelStrong, message-carrying) ---
\node[labelStrong, text=cGray!85!black] at (2.30, 2.62) {$I(t) \sim t^{-n}$};

% --- Axis labels (minimal, no tick numbers) ---
\node[labelStd, text=cGray, anchor=west]             at (4.32, 0.45) {$\log\,t$};
\node[labelStd, text=cGray, rotate=90, anchor=south] at (0.10, 1.45) {$\log\,I$};

% ============================================================
% Panel E — V_s(t) ISPD raw decay   x: 4.50..6.95, y: 0.10..4.45
% Single log-log plot, schematic open-circle data — paired with Panel F.
% ============================================================

\draw[cGray, line width=0.40pt, -{Stealth[length=2.5pt,width=1.8pt]}]
  (4.70, 0.85) -- (6.85, 0.85);
\draw[cGray, line width=0.40pt, -{Stealth[length=2.5pt,width=1.8pt]}]
  (4.70, 0.85) -- (4.70, 3.95);

% Power-law decay envelope (log-log → straight line w/ slope)
\draw[cRed!75!black, line width=0.75pt, dashed]
  (4.85, 3.65) -- (6.75, 1.20);

% Open-circle data markers (schematic, ~7 points along envelope)
\foreach \px/\py in {%
  4.95/3.55, 5.25/3.20, 5.55/2.85, 5.85/2.50,
  6.15/2.15, 6.45/1.80, 6.75/1.45%
} {
  \fill[white] (\px, \py) circle (0.95mm);
  \draw[cRed!75!black, line width=0.45pt] (\px, \py) circle (0.95mm);
}

% Axis labels
\node[labelStd, text=cGray, anchor=west] at (6.90, 0.85) {$\log\,t$};
\node[labelStd, text=cGray, rotate=90, anchor=south] at (4.45, 2.40) {$\log V_s$};

% Subtitle (labelMute)
\node[labelMute] at (5.78, 4.18) {surface potential decay};

% ============================================================
% E ↔ F paired connector — "ISPD" transformation
% Short arrow + label between Panel E (right) and Panel F (left).
% ============================================================
\draw[cGray!65!black, line width=0.55pt,
      -{Stealth[length=3.5pt,width=2.5pt]}]
  (6.88, 2.40) -- (7.17, 2.40);
\node[labelMute, anchor=south] at (7.02, 2.50) {ISPD};

% ============================================================
% Panel F — g(E_t) bimodal trap DOS  (HERO #2)   x: 7.00..10.45, y: 0.10..4.45
% Deep peak height = shallow × 1.8  (design.md §4 / Q4)
% Plot frame style: shared with Panel E (paired ISPD) — defined inline here,
% Panel E reuses same axis weights/fonts when implemented.
% ============================================================

% --- Plot axes (cGray, thin Stealth) ---
\draw[cGray, line width=0.40pt, -{Stealth[length=2.5pt,width=1.8pt]}]
  (7.20, 0.85) -- (10.40, 0.85);
\draw[cGray, line width=0.40pt, -{Stealth[length=2.5pt,width=1.8pt]}]
  (7.20, 0.85) -- (7.20, 3.95);

% --- Shallow Gaussian (left, blue, peak height 1.50) ---
\BellCurve[draw=cBlue!75!black, fill=cBlue!22, line width=0.65pt]%
  {7.65, 0.85, 8.55, 2.35, up}

% --- Deep Gaussian (right, red, peak height 2.70 ≈ 1.8× shallow) ---
\BellCurve[draw=cRed!80!black, fill=cRed!28, line width=0.75pt]%
  {9.00, 0.85, 9.90, 3.55, up}

% --- τ_d double-headed arrow (between Shallow and Deep peaks, briefing §2/§3) ---
% C005/FN-1 fix: τ_d arrow was missing entirely from .tex despite briefing requirement.
\draw[cRed!80!black, line width=0.85pt, dashed,
      {Stealth[length=4pt,width=3pt]}-{Stealth[length=4pt,width=3pt]}]
  (8.55, 3.75) -- (9.90, 3.75);
\node[labelStrong, text=cRed!80!black] at (9.22, 3.92) {$\tau_d$};

% --- Axis labels (only x-axis label; y-axis sense carried by subtitle) ---
\node[labelStd, text=cGray, anchor=west] at (10.45, 0.85) {$E_t$};

% --- Region labels under x-axis (labelStrong tier, color-linked to Panel C) ---
\node[labelStrong, text=cBlue!75!black] at (8.10, 0.55) {Shallow};
\node[labelStrong, text=cRed!80!black]  at (9.45, 0.55) {Deep};

% --- Panel F subtitle (labelMute) ---
\node[labelMute] at (8.75, 4.18) {trap DOS (ISPD-derived)};

% ============================================================
% Panel G — Mechanical evidence (bent cantilever, v7-pattern)
% Re-authored 2026-05-08 v3: adopt sc7_concept_tikz_v7.tex pattern for natural
% cantilever bending — single control-point quadratic bezier + two-stroke
% layering (dark outer + light inner) + line cap=round + dashed reference line.
% Layout per briefing §3 Panel G: electrode LEFT, polymer strip RIGHT (bending
% AWAY from electrode). Panel area x:10.50..13.95, y:0.10..4.45.
% ============================================================

% --- Electrode (LEFT, flat dark gray with subtle horizontal gradient) ---
\path[draw=cGray!75!black, line width=0.55pt, shade,
      left color=cGray!75!black, right color=cGray!50!black]
  (10.95, 0.50) rectangle (11.40, 3.80);

% --- Top clamp / fixture (anchor for cantilever, dark slab) ---
\filldraw[fill=cGray!55!black, draw=cGray!75!black, line width=0.30pt,
          rounded corners=0.4mm]
  (12.30, 3.78) rectangle (12.70, 3.95);

% --- Vertical reference line (un-bent strip silhouette, dashed faint) ---
% Visualizes the bending displacement: tip deflects rightward from this line.
\draw[cGray!35, line width=0.35pt, dashed] (12.50, 0.70) -- (12.50, 3.78);

% --- Cantilever strip — v7 two-stroke bezier (single control point) ---
% Outer stroke: darker amber, thicker (6pt = ~0.21cm visual width).
% Inner stroke: lighter amber, thinner (4pt). Edge-highlight pseudo-3D effect.
% Single control point at (12.65, 1.95) gives soft natural bend, tip displaced
% ~0.7cm rightward (cantilever deflection visualization).
\draw[cAmber!85!black, line width=6pt, line cap=round]
  (12.50, 3.78) ..controls(12.65, 1.95).. (13.20, 0.85);
\draw[cAmber!55!white, line width=4pt, line cap=round]
  (12.50, 3.78) ..controls(12.65, 1.95).. (13.20, 0.85);

% --- q_tr trapped charges along strip's electrode-facing (LEFT) edge ---
% Positions computed at bezier params t≈0.20, 0.45, 0.70, 0.92 with
% perpendicular offset toward smaller x (~0.10cm to land on left edge).
\foreach \qx/\qy in {%
  12.50/3.05,
  12.62/2.30,
  12.80/1.55,
  13.00/0.95%
} {
  \node[circle, draw=cBlue!85!black, fill=cBlue!65,
        inner sep=0pt, minimum size=1.6mm, line width=0.30pt,
        font=\sffamily\fontsize{5}{6}\selectfont, text=white]
    at (\qx, \qy) {$-$};
}

% --- q_tr label (top of charge stack) ---
\node[labelStd, text=cBlue!75!black, anchor=south east]
  at (12.42, 3.20) {$q_{tr}$};

% --- Air gap label (between electrode and strip mid-region) ---
\node[labelMute] at (11.85, 1.95) {air gap};

% --- Coulomb repulsion (red, polymer pushed AWAY from electrode = RIGHTWARD) ---
% Tail at strip mid-right, head outside strip to the right.
\draw[cRed!80!black, line width=1.1pt, -{Stealth[length=5.5pt,width=4.5pt]}]
  (13.05, 2.55) -- (13.65, 2.55);
\node[labelStrong, text=cRed!80!black, anchor=west]
  at (13.05, 2.85) {Coulomb};
\node[labelStd, text=cRed!75!black, anchor=west]
  at (13.05, 2.30) {repulsion};

% --- Maxwell attraction (blue, polymer pulled TOWARD electrode = LEFTWARD) ---
% Tail at strip lower-left, head pointing into air gap.
\draw[cBlue!75!black, line width=1.1pt, -{Stealth[length=5.5pt,width=4.5pt]}]
  (12.55, 1.20) -- (11.95, 1.20);
\node[labelStrong, text=cBlue!75!black, anchor=west]
  at (12.60, 1.45) {Maxwell};
\node[labelStd, text=cBlue!70!black, anchor=west]
  at (12.60, 0.95) {attraction};

% ============================================================
% Inter-row connector — "convergent evidence"
% Drops at x=10.42 ≈ F-G boundary (Row 2 visual midpoint).
% ============================================================
% C006 fix: "convergent evidence" label removed (not in briefing).
% Connector arrow alone carries the Row 1 → Row 2 transition.
\draw[interArrow] (10.42, 5.00) -- (10.42, 4.55);

\end{tikzpicture}%
}
\end{document}
